\documentclass[12pt,a4paper]{article}
\usepackage{lmodern}
\usepackage[T1]{fontenc}
\usepackage{amsmath,amssymb,amsthm}
\usepackage[margin=1.2in]{geometry}
\usepackage[dvipsnames,svgnames,x11names]{xcolor}
\usepackage{hyperref}
\usepackage{framed}
\usepackage{enumitem}
\usepackage{mdframed}
\usepackage{array}
\usepackage{booktabs}

\hypersetup{colorlinks=true,linkcolor=NavyBlue,urlcolor=NavyBlue}

% Define named colors for boxes
\definecolor{cleanbg}{RGB}{235,248,235}
\definecolor{cleanborder}{RGB}{50,140,50}
\definecolor{openbg}{RGB}{252,235,235}
\definecolor{openborder}{RGB}{180,50,50}
\definecolor{notebg}{RGB}{235,242,252}
\definecolor{noteborder}{RGB}{60,100,180}

% Clean / Open status boxes using mdframed
\newmdenv[
  backgroundcolor=cleanbg,
  linecolor=cleanborder,
  linewidth=1.5pt,
  innertopmargin=6pt,
  innerbottommargin=6pt,
  skipabove=8pt,
  skipbelow=8pt
]{cleanenv}

\newmdenv[
  backgroundcolor=openbg,
  linecolor=openborder,
  linewidth=1.5pt,
  innertopmargin=6pt,
  innerbottommargin=6pt,
  skipabove=8pt,
  skipbelow=8pt
]{openenv}

\newmdenv[
  backgroundcolor=notebg,
  linecolor=noteborder,
  linewidth=1pt,
  innertopmargin=6pt,
  innerbottommargin=6pt,
  skipabove=8pt,
  skipbelow=8pt
]{notebox}

\newtheorem{theorem}{Theorem}
\newtheorem{lemma}[theorem]{Lemma}
\newtheorem{proposition}[theorem]{Proposition}
\newtheorem{corollary}[theorem]{Corollary}
\newtheorem{definition}{Definition}
\newtheorem{remark}{Remark}
\theoremstyle{definition}
\newtheorem{subproblem}{Subproblem}

\title{\textbf{Cancelling Translated Copies of Points in Convex Position}\\[6pt]
\large Progress Report}
\author{}
\date{May 2026}

\begin{document}
\maketitle
\tableofcontents
\newpage

%%%%%%%%%%%%%%%%%%%%%%%%%%%%%%%%%%%%%%%%%%%%%%%%%%%%%%%%%%%%%
\section{The Problem}
%%%%%%%%%%%%%%%%%%%%%%%%%%%%%%%%%%%%%%%%%%%%%%%%%%%%%%%%%%%%%

\begin{definition}
A finite set $S$ of points in $\mathbb{R}^2$ is in \emph{convex position} if every point of $S$
is a vertex of its convex hull $\operatorname{conv}(S)$.
\end{definition}

\begin{definition}[Setup]
Let $S$ be a finite set of $n \geq 1$ points in $\mathbb{R}^2$ in convex position.
Let $T$ be a finite set of vectors in $\mathbb{R}^2$, each assigned a sign $\sigma(t)\in\{+1,-1\}$.
For every $p\in\mathbb{R}^2$ define
\[
  \sigma(p) \;=\; \sum_{\substack{t\in T\\ p-t\in S}} \sigma(t).
\]
Write $\operatorname{supp}(\sigma):=\{p\in\mathbb{R}^2:\sigma(p)\neq 0\}$, and set
$T_\pm:=\{t\in T:\sigma(t)=\pm1\}$.
\end{definition}

\begin{theorem}[Lower Bound --- proved]
$|\operatorname{supp}(\sigma)|\ge n$.
\end{theorem}

\begin{theorem}[Main Theorem --- partially proved]\label{thm:main}
Assume $n\ge 3$ and $T$ contains at least one positive and one negative vector.
If $|\operatorname{supp}(\sigma)|=n$ then $n=4$ and $S$ is the vertex set of a parallelogram.
\end{theorem}

\begin{remark}
The condition ``both signs present'' is necessary: a single translate gives $|\operatorname{supp}(\sigma)|=n$
for any $S$.  The theorem says: once genuine cancellation occurs but the support is still minimal,
$S$ is forced to be a parallelogram (and in particular $n=4$).
\end{remark}

%%%%%%%%%%%%%%%%%%%%%%%%%%%%%%%%%%%%%%%%%%%%%%%%%%%%%%%%%%%%%
\section{Proof Strategy}
%%%%%%%%%%%%%%%%%%%%%%%%%%%%%%%%%%%%%%%%%%%%%%%%%%%%%%%%%%%%%

The proof is decomposed into six subproblems.  Write $P:=\operatorname{conv}(S)$,
$Q:=\operatorname{conv}(T)$.  Label the vertices of $P$ as $s_1,\ldots,s_n$
counterclockwise (indices mod $n$), with edge vectors $a_i:=s_{i+1}-s_i$.

\medskip
\begin{center}
\renewcommand{\arraystretch}{1.4}
\begin{tabular}{clc}
\toprule
\textbf{SP} & \textbf{Content} & \textbf{Status}\\
\midrule
1 & Lonely Points Lemma + equality rigidity & \textcolor{green!60!black}{\textbf{clean}}\\
2 & Exposed-translate theorem & \textcolor{green!60!black}{\textbf{clean}}\\
3 & Boundary Minkowski structure & \textcolor{green!60!black}{\textbf{clean}}\\
4 & Alternating-chain structure on each face & \textcolor{green!60!black}{\textbf{clean}}\\
5 & No three independent edge directions & \textcolor{red!70!black}{\textbf{open}}\\
6 & Two-direction case forces parallelogram & \textcolor{green!60!black}{\textbf{clean}}\\
\bottomrule
\end{tabular}
\end{center}

\bigskip
\noindent\textbf{How they combine.}
Subproblems 1--4 establish the rigid algebraic--geometric structure of any equality
configuration.  Subproblem 5 rules out three or more independent edge directions
(the hard case).  Subproblem 6 shows the remaining two-direction case forces
$n=4$ and a parallelogram.  Everything is proved \emph{except} Subproblem 5.

%%%%%%%%%%%%%%%%%%%%%%%%%%%%%%%%%%%%%%%%%%%%%%%%%%%%%%%%%%%%%
\section{Subproblem 1: Lonely Points Lemma}
%%%%%%%%%%%%%%%%%%%%%%%%%%%%%%%%%%%%%%%%%%%%%%%%%%%%%%%%%%%%%

\begin{cleanenv}
\textbf{Status: VERIFIED CLEAN} --- adversarial checker found no errors (3/3 independent runs).
\end{cleanenv}

\begin{subproblem}[Lonely Points + Equality Rigidity]\label{sp:1}
Under the general setup (no equality hypothesis yet):
\begin{enumerate}[label=(\alph*)]
\item \textbf{Lower bound:} $|\operatorname{supp}(\sigma)|\ge n$.
\item \textbf{Equality rigidity:} If $|\operatorname{supp}(\sigma)|=n$, then:
  \begin{enumerate}[label=(\roman*)]
  \item For each $s\in S$ there is a \emph{unique} $t_s\in T$ such that
        $\ell_s:=s+t_s$ is lonely (has a unique representation in $S+T$).
  \item The lonely points $\ell_s$ are pairwise distinct and
        $\operatorname{supp}(\sigma)=\{\ell_s:s\in S\}$.
  \item $\sigma(\ell_s)=\sigma(t_s)\in\{\pm1\}$.
  \item $\bigl||T_+|-|T_-|\bigr|\le 1$.
  \end{enumerate}
\end{enumerate}
\end{subproblem}

\subsection*{Proof sketch}

A point $p\in S+T$ is \emph{lonely} if it has a unique representation $p=s+t$.
For each $s\in S$, find a vector $u_s$ that strictly exposes $s$ in $S$
(i.e.\ $\langle u_s,s\rangle > \langle u_s,s'\rangle$ for all $s'\ne s$; constructed by
taking the nearest point in $\operatorname{conv}(S\setminus\{s\})$).
Let $t_s$ maximise $\langle u_s,\cdot\rangle$ on $T$.
Then $\ell_s:=s+t_s$ is lonely: any other representation $s'+t'=\ell_s$
would require $\langle u_s,s'\rangle+\langle u_s,t'\rangle=\langle u_s,s\rangle+\langle u_s,t_s\rangle$,
forcing $s'=s$ and $t'=t_s$ by strict maximality.
Since $\sigma(\ell_s)=\sigma(t_s)\ne 0$, we get $n$ distinct support points.

For equality rigidity: if $s$ had two lonely points $s+t$ and $s+t'$ with $t\ne t'$,
we get $(n-1)+2=n+1$ distinct support points, a contradiction.
The balance estimate follows from $\sum_p\sigma(p)=n(|T_+|-|T_-|)$
combined with $|\sum_p\sigma(p)|\le n$. $\square$

%%%%%%%%%%%%%%%%%%%%%%%%%%%%%%%%%%%%%%%%%%%%%%%%%%%%%%%%%%%%%
\section{Subproblem 2: Exposed-Translate Theorem}
%%%%%%%%%%%%%%%%%%%%%%%%%%%%%%%%%%%%%%%%%%%%%%%%%%%%%%%%%%%%%

\begin{cleanenv}
\textbf{Status: VERIFIED CLEAN} --- adversarial checker found no errors.
\end{cleanenv}

Under the equality hypothesis, write $t_i:=t_{s_i}$ for the unique translate from SP~1
attached to vertex $s_i$.  Define the normal cone
\[
  C_i := \{u\in\mathbb{R}^2 : \text{$s_i$ is the unique maximiser of $\langle u,\cdot\rangle$ on $S$}\}.
\]

\begin{subproblem}[Exposed-Translate Theorem]\label{sp:2}
For every $u\in C_i$, $t_i$ is the unique maximiser of $\langle u,\cdot\rangle$ on $T$.
\end{subproblem}

\subsection*{Proof sketch}

If $u\in C_i$ and $t\in T$ maximises $\langle u,\cdot\rangle$ on $T$, the same exposed-maximum
argument as in SP~1 shows $s_i+t$ is lonely.
By SP~1 uniqueness, $t=t_i$.
So $\operatorname{argmax}_T\langle u,\cdot\rangle=\{t_i\}$ for every $u\in C_i$. $\square$

%%%%%%%%%%%%%%%%%%%%%%%%%%%%%%%%%%%%%%%%%%%%%%%%%%%%%%%%%%%%%
\section{Subproblem 3: Boundary Minkowski Structure}
%%%%%%%%%%%%%%%%%%%%%%%%%%%%%%%%%%%%%%%%%%%%%%%%%%%%%%%%%%%%%

\begin{cleanenv}
\textbf{Status: VERIFIED CLEAN} --- adversarial checker found no errors (3 runs).
\end{cleanenv}

\begin{subproblem}[Boundary Minkowski Structure]\label{sp:3}
Under the equality hypothesis:
\begin{enumerate}[label=(\arabic*)]
\item The points $x_i:=s_i+t_i$ are exactly the vertices of $R:=P+Q$, in cyclic order.
\item For each $i$, the exposed face of $Q$ in the direction normal to the edge $[s_i,s_{i+1}]$
      of $P$ is the segment $[t_i,t_{i+1}]$, and
      \[
        t_{i+1}-t_i = \lambda_i\,a_i \quad\text{for some }\lambda_i\ge 0.
      \]
\end{enumerate}
\end{subproblem}

\subsection*{Proof sketch}

Let $u_i$ be an outer normal to $[s_i,s_{i+1}]$.  Small perturbations
$u_i\pm\varepsilon a_i$ (for small $\varepsilon>0$) enter cones $C_i$ and $C_{i+1}$
respectively; by SP~2, $t_i$ and $t_{i+1}$ are the unique maximisers of those
perturbed directions on $T$.  Letting $\varepsilon\to 0$, both $t_i$ and $t_{i+1}$
maximise $\langle u_i,\cdot\rangle$ on $T$, so both lie in the exposed face
$G_i:=F_Q(u_i)$.  Since $G_i$ is a compact convex segment parallel to $a_i$,
and $t_i$, $t_{i+1}$ are respectively the minimum and maximum of $\langle a_i,\cdot\rangle$
on $G_i$, we get $G_i=[t_i,t_{i+1}]$ by the interval structure of compact convex
subsets of a line.  That $\lambda_i\ge 0$ follows because $x_i$ must be an endpoint
(not interior) of the corresponding edge of $R$.

For part (1): the Minkowski-sum face formula $F_{P+Q}(u)=F_P(u)+F_Q(u)$
gives $F_R(u)=\{x_i\}$ for $u\in C_i$ (a vertex), and $F_R(u_i)=[x_i,x_{i+1}]$
(an edge).  A cyclic-normal-order lemma (exterior angles sum to $2\pi$) shows these
are exactly all vertices and edges of $R$. $\square$

%%%%%%%%%%%%%%%%%%%%%%%%%%%%%%%%%%%%%%%%%%%%%%%%%%%%%%%%%%%%%
\section{Subproblem 4: Alternating Chains on Faces}
%%%%%%%%%%%%%%%%%%%%%%%%%%%%%%%%%%%%%%%%%%%%%%%%%%%%%%%%%%%%%

\begin{cleanenv}
\textbf{Status: VERIFIED CLEAN} --- adversarial checker found no errors.
\end{cleanenv}

For each $i$, let $F_i:=[t_i,t_{i+1}]\cap T$ be the translates on the $i$-th face.

\begin{subproblem}[Alternating-Chain Structure]\label{sp:4}
Under the equality hypothesis, for each $i$:
\begin{enumerate}[label=(\arabic*)]
\item There exists a non-negative integer $m_i$ such that $t_{i+1}=t_i+m_i\,a_i$ and
\[
  F_i = \{t_i,\; t_i+a_i,\; t_i+2a_i,\;\ldots,\; t_i+m_i\,a_i\}.
\]
\item The signs alternate along the chain:
      $\sigma(t_i+k\,a_i)=(-1)^k\,\sigma(t_i)$ for $0\le k\le m_i$.
\end{enumerate}
\end{subproblem}

\subsection*{Proof sketch}

Any interior point of edge $[x_i,x_{i+1}]$ of $R$ has $\sigma=0$
(by SP~1, the support equals $\{x_1,\ldots,x_n\}$, the vertex set of $R$).
For $t=t_i+\lambda a_i\in F_i$ with $\lambda<m_i$, the point $p=s_{i+1}+t$
is interior to that edge and must cancel.  The convex-position of $S$ forces
the only representations of $p$ in $S+T$ to use translates $t$ or $t+a_i$;
since $\sigma(t)\ne 0$, the translate $t+a_i$ must also lie in $T$ with
opposite sign.  This gives a successor rule $\lambda\in\Lambda\Rightarrow\lambda+1\in\Lambda$,
and a predecessor rule $\lambda>0\Rightarrow\lambda-1\in\Lambda$.
No fractional parameter can lie in $\Lambda$ (it would give a lonely interior
point with $\sigma\ne 0$), so $\Lambda\subset\mathbb{Z}$, and we conclude
$\Lambda=\{0,1,\ldots,m_i\}$ with alternating signs. $\square$

%%%%%%%%%%%%%%%%%%%%%%%%%%%%%%%%%%%%%%%%%%%%%%%%%%%%%%%%%%%%%
\section{Subproblem 5: No Three Independent Directions}
%%%%%%%%%%%%%%%%%%%%%%%%%%%%%%%%%%%%%%%%%%%%%%%%%%%%%%%%%%%%%

\begin{openenv}
\textbf{Status: OPEN} --- This is the remaining hard case.  Multiple attempts
(versions v1--v19) have all contained critical errors.  The problem is unsolved.
\end{openenv}

\begin{subproblem}[No Three Independent Directions]\label{sp:5}
Under the equality hypothesis with the structure of SP~1--4, the edge vectors
$a_1,\ldots,a_n$ of $P$ span at most two unoriented directions
(i.e.\ every $a_i$ is a real multiple of one of two fixed vectors).
\end{subproblem}

\subsection*{What is known}

The following are proved and available.

\begin{enumerate}
\item \textbf{L$^2$-norm identity.}  Since $\sigma(\ell_s)^2=1$:
\[
  n(|T|-1) \;=\; \underbrace{\sum_{\substack{t_1\ne t_2\\\sigma(t_1)\sigma(t_2)=-1}}
  |S\cap(S+(t_1-t_2))|}_{\text{opp-sign sum}}
  \;-\;
  \underbrace{\sum_{\substack{t_1\ne t_2\\\sigma(t_1)\sigma(t_2)=+1}}
  |S\cap(S+(t_1-t_2))|}_{\text{same-sign sum}}.
\]

\item \textbf{Overlap bound.}  For any $v\ne 0$, $|S\cap(S+v)|\le 2$
(since $S$ is in convex position).

\item \textbf{Coarse L$^2$ bound.}  From (1) and (2):
\[
  n(|T|-1) \;\le\; 4\,|T_+|\,|T_-|.
\]
This gives $n\le 4$ when $|T_+|=1$ or $|T_-|=1$, but \emph{fails} for
$|T_+|=|T_-|=k\ge 2$.

\item \textbf{Alternating chains} (SP~4): each face of $P$ carries an arithmetic
progression of translates with alternating signs.
\end{enumerate}

\subsection*{Known dead-ends}

\begin{itemize}
\item Assuming $|T|=2$ is unjustified; explicit equality configurations with $|T_+|=|T_-|=2$ exist (e.g.\ $S$ a parallelogram, $T=\{0,a+b,-a,-b\}$).
\item The support need not equal the vertex set of $\operatorname{conv}(S+T)$.
\item Polynomial/lattice arguments do not apply in $\mathbb{R}^2$.
\item No ``opposite-sign lonely point'' is guaranteed when all lonely signs are equal.
\end{itemize}

\subsection*{Most promising open approach}

Lower-bound the same-sign sum to tighten the L$^2$ estimate.
For $|T_+|=|T_-|=k\ge 2$ there are $2k(k-1)$ same-sign ordered pairs.
If each contributes $\ge 1$ to the same-sign sum, then
\[
  n(2k-1) \;\le\; 4k^2 - 2k(k-1) \;=\; 2k^2+2k,
\]
giving $n\le 4$ for all $k\ge 2$.  The needed ingredient is a proof that
$|S\cap(S+(t_1-t_2))|\ge 1$ for same-sign pairs $(t_1,t_2)$, or an alternative
argument using the alternating-chain structure from SP~4.

%%%%%%%%%%%%%%%%%%%%%%%%%%%%%%%%%%%%%%%%%%%%%%%%%%%%%%%%%%%%%
\section{Subproblem 6: Two Directions Forces a Parallelogram}
%%%%%%%%%%%%%%%%%%%%%%%%%%%%%%%%%%%%%%%%%%%%%%%%%%%%%%%%%%%%%

\begin{cleanenv}
\textbf{Status: VERIFIED CLEAN} --- adversarial checker found no errors (3/3 independent runs).
\end{cleanenv}

\begin{subproblem}[Two Directions Forces a Parallelogram]\label{sp:6}
If $S$ is in convex position with $n\ge 3$ and the edge vectors of $\operatorname{conv}(S)$
span at most two unoriented directions, then $n=4$ and $S$ is the vertex set of a parallelogram.
\end{subproblem}

\subsection*{Proof}

\textbf{Step 1: Adjacent edges are non-parallel.}
If $e_i\parallel e_{i+1}$, then $v_i,v_{i+1},v_{i+2}$ are collinear, one being a strict
convex combination of the other two.  This contradicts convex position.
So the two direction classes alternate around the polygon: $[u],[w],[u],[w],\ldots$,
forcing $n$ even.

\textbf{Step 2: $n\le 4$.}
The exterior angles $\varepsilon_i>0$ sum to $2\pi$, so the $n$ directed edge directions
$e_1,\ldots,e_n$ are pairwise distinct (their lifted arguments
$\theta_1<\cdots<\theta_n<\theta_1+2\pi$ are all in the same period).
Since there are only two unoriented directions, there are at most four possible
oriented directions $u,w,-u,-w$.  Hence $n\le 4$.

\textbf{Step 3: $n=4$, parallelogram.}
With $n$ even, $n\ge 3$, $n\le 4$, we get $n=4$.
Write the four edges as $B-A=au$, $C-B=bw$, $D-C=\sigma cu$, $A-D=\tau dw$
(with $a,b,c,d>0$, $\sigma,\tau\in\{\pm1\}$).
The closure condition $(a+\sigma c)u+(b+\tau d)w=0$ and linear independence of
$u,w$ force $\sigma=\tau=-1$, $a=c$, $b=d$.
So opposite sides are parallel and equal: $ABCD$ is a parallelogram. $\square$

%%%%%%%%%%%%%%%%%%%%%%%%%%%%%%%%%%%%%%%%%%%%%%%%%%%%%%%%%%%%%
\section{Summary}
%%%%%%%%%%%%%%%%%%%%%%%%%%%%%%%%%%%%%%%%%%%%%%%%%%%%%%%%%%%%%

The main theorem reduces entirely to Subproblem~5.  All other ingredients are proved
and verified.

\begin{theorem}[Conditional]\label{thm:cond}
If Subproblem~\ref{sp:5} holds, then Theorem~\ref{thm:main} holds in full:
every equality configuration with $n\ge 3$ and both signs in $T$ forces $n=4$
and $S$ a parallelogram.
\end{theorem}

\begin{proof}
SP~5 provides at most two edge directions; SP~6 then gives $n=4$ and a parallelogram.
\end{proof}

\bigskip
\begin{notebox}
\textbf{What remains.}
\begin{itemize}
\item Prove Subproblem~5.
\item The key gap is the case $|T_+|=|T_-|=k\ge 2$;
      the L$^2$ bound alone is insufficient here.
\item A lower bound $|S\cap(S+(t_1-t_2))|\ge 1$ for same-sign pairs,
      combined with the L$^2$ identity, would complete the proof.
\item See \texttt{proof/feedback\_round5.txt} for a full account of failed approaches
      and a precise statement of the key missing lemma.
\end{itemize}
\end{notebox}

\bigskip
\noindent\textbf{File guide.}

\medskip
\noindent
\begin{tabular}{@{}ll@{}}
\texttt{proof/section\_1\_lonely\_rigidity\_sol\_v3.tex} & SP~1 (clean)\\
\texttt{proof/section\_2\_exposed\_translate\_sol\_v1.tex} & SP~2 (clean)\\
\texttt{proof/section\_3\_boundary\_minkowski\_sol\_v5.tex} & SP~3 (clean)\\
\texttt{proof/section\_4\_alternating\_chains\_sol\_v2.tex} & SP~4 (clean)\\
\texttt{proof/section\_5\_three\_directions.tex} & SP~5 problem statement\\
\texttt{proof/feedback\_round5.txt} & SP~5 failure diagnosis\\
\texttt{proof/section\_6\_two\_directions\_sol\_v3.tex} & SP~6 (clean)\\
\end{tabular}

\end{document}
